\documentclass[10pt,a4paper]{amsart}
\usepackage[margin=19mm]{geometry}
\usepackage{amsmath,amssymb,mathtools}
\usepackage[scr=rsfs,cal=euler]{mathalfa}
\usepackage{xfrac}
\usepackage[unicode,hidelinks,bookmarks=false]{hyperref}
\newcommand{\ie}{i.e.\ }
\newcommand{\cf}{cf.\ }
\newcommand{\resp}{respectively\ }
\newcommand{\Subsection}{\subsection}
\providecommand{\nobreakdash}{\nobreak\hbox{-}}
\setlength{\emergencystretch}{2em}
\multlinegap0pt
\usepackage{amssymb}
\newtheorem{prop}{Proposition}
\def\Cal{\mathcal}
\newcommand{\Ga}{\Gamma}
\newcommand{\La}{\Lambda}
\newcommand{\Ph}{\Phi}
\newcommand{\Ric}{\operatorname{Ric}}
\def\frak{\mathfrak}
\newcommand{\la}{\lambda}
\newcommand{\si}{\sigma}
\newcommand{\tn}{\widetilde{\nabla}}
\title{A new construction of the\\ Riemannian deformation sequence}
\author{Andreas Cap}
\date{Source: arXiv:2607.26653v1 (CC BY 4.0)}
\begin{document}
\maketitle

\noindent\emph{Source and licence.} A new construction of the\\ Riemannian deformation sequence. Version arXiv:2607.26653v1 (CC BY 4.0), distributed by arXiv under the Creative Commons Attribution 4.0 International licence, http://creativecommons.org/licenses/by/4.0/, as recorded in the arXiv licence metadata for this exact version and shown on the abstract page https://arxiv.org/abs/2607.26653v1. Full licence notice: \texttt{licenses/arxiv-2607.26653-CC-BY-4.0.txt}.\par\smallskip

\noindent\textit{Editorial scope.} Counted unit: the article's prop prop2.2 (source0.tex lines 417--432). The article's complete original proof of it is delivered with it (source0.tex lines 433--523, one \texttt{proof} environment of 91 lines). No mathematics is written, repaired or re-derived: the statement and the proof are the article's own lines, in the article's own notation, and no retained line is substituted or re-flowed.  Retained uncounted as the source-local context the proof needs: lines 371--374; lines 378--380; lines 385--387; lines 407--412.  Nothing was omitted from any retained span, so the package carries no omission record.  No retained line contains a citation, so the wrapper carries no bibliography and no bibliography entry.  No retained line contains a figure environment, an image-inclusion command or drawing code, so no image is delivered and the package declares no asset dependency.  Editorial formatting only, in addition: an AMS \texttt{amsart} preamble that restates the article's own declarations and macros used by the retained lines, namely the environments \texttt{prop} on separate counters, together with the standard packages those lines need.  The retained environments print in this document as Proposition 1 in order of appearance, whereas the article prints the same items with the numbering of its own full document.  The complete native source of the article as distributed in the arXiv e-print for this exact version is bundled unchanged as \texttt{sources/206284/source0.tex}.
\begin{equation}\label{curv-W}
  \nabla_{\xi_1}\nabla_{\xi_2}\si-\nabla_{\xi_2}\nabla_{\xi_1}\si-\nabla_{[\xi_1,\xi_2]}\si=
  R(\xi_1,\xi_2)\bullet\si.
\end{equation}

\begin{equation}\label{Bianchi1}
R(\xi_1,\xi_2)(\xi_3)+R(\xi_3,\xi_1)(\xi_2)+R(\xi_2,\xi_3)(\xi_1)=0. 
\end{equation}

\begin{equation}\label{Bianchi2}
(\nabla_{\xi_1}R)(\xi_2,\xi_3)+(\nabla_{\xi_3}R)(\xi_1,\xi_2)+(\nabla_{\xi_2}R)(\xi_3,\xi_1)=0.
\end{equation}

\begin{equation}\label{conn-def}
\nabla^{\Cal A}_\xi\begin{pmatrix} \eta\\ \Phi \end{pmatrix}:=\begin{pmatrix}
\nabla_\xi\eta-\Phi(\xi)\\ \nabla_\xi\Phi \end{pmatrix} \qquad 
\tn^{\Cal A}_\xi\begin{pmatrix} \eta\\ \Phi \end{pmatrix}:=\begin{pmatrix}
\nabla_\xi\eta-\Phi(\xi)\\ \nabla_\xi\Phi-R(\xi,\eta) \end{pmatrix}
\end{equation}

\begin{prop}\label{prop2.2}
Both $\nabla^{\Cal A}$ and $\tn^{\Cal A}$ define linear connections on the bundle
$\Cal AM$ that are canonically associated to $(M,g)$. Denoting their curvatures by
$R^{\Cal A}$ and $\widetilde{R}^{\Cal A}$, respectively, we obtain
$$
R^{\Cal A}(\xi_1,\xi_2)\begin{pmatrix} \eta\\ \Phi \end{pmatrix}=\begin{pmatrix}
R(\xi_1,\xi_2)(\eta) \\ R(\xi_1,\xi_2)\bullet \Phi \end{pmatrix}
$$

$$ \widetilde{R}^{\Cal A}(\xi_1,\xi_2)\begin{pmatrix}
  \eta\\ \Phi \end{pmatrix}=\begin{pmatrix} 0 \\ (\nabla_\eta
R)(\xi_1,\xi_2)-(\Phi\bullet R)(\xi_1,\xi_2)\end{pmatrix}
$$
In particular, the connection $\nabla^{\Cal A}$ is flat if and only if $g$ is flat,
while $\tn^{\Cal A}$ is flat if and only if $g$ has constant sectional curvature.
\end{prop}

\begin{proof}
Both expressions in \eqref{conn-def} are the sum of the component-wise Levi-Civita
connection and a tensorial term, so they both define linear connections on $\Cal
AM$. To compute the curvatures we directly evaluate the defining equation for the
curvature of a linear connection. Starting with $\nabla^{\Cal A}$, differentiating the
first formula in \eqref{conn-def} shows that we obtain for $\xi_1,\xi_2,\eta\in\frak
X(M)$ and $\Ph\in\frak o(TM)$ the equation
\begin{equation}\label{n2}
\nabla^{\Cal A}_{\xi_1}\nabla^{\Cal
  A}_{\xi_2}\begin{pmatrix}\eta\\ \Phi \end{pmatrix}=
\begin{pmatrix}
  \nabla_{\xi_1}\nabla_{\xi_2}\eta-\nabla_{\xi_1}\Phi(\xi_2)-(\nabla_{\xi_2}\Phi)(\xi_1)\\
  \nabla_{\xi_1}\nabla_{\xi_2}\Phi
\end{pmatrix}. 
\end{equation}
To obtain $R^{\Cal A}(\xi_1,\xi_2)$ we have to subtract from this the same expression
with $\xi_1$ and $\xi_2$ exchanged as well as
\begin{equation}\label{brack-term}
  \begin{pmatrix} \nabla_{[\xi_1,\xi_2]}\eta-\Ph([\xi_1,\xi_2])\\
    \nabla_{[\xi_1,\xi_2]}\Phi\end{pmatrix}.
\end{equation}    
In the $\frak o(TM)$-component, this readily produces $R(\xi_1,\xi_2)\bullet \Phi$ by
\eqref{curv-W}. In the $TM$-component, we obtain the sum of $R(\xi_1,\xi_2)(\eta)$
and
\begin{equation}\label{nb}
-\nabla_{\xi_1}\Phi(\xi_2)-(\nabla_{\xi_2}\Phi)(\xi_1)+\nabla_{\xi_2}\Phi(\xi_1)+
(\nabla_{\xi_1}\Phi)(\xi_2)-\Phi([\xi_1,\xi_2]). 
\end{equation}
Now the first and fourth term add up to $\Ph(\nabla_{\xi_1}\xi_2)$, while the second
and third term add up to $-\Ph(\nabla_{\xi_2}\xi_1)$. But torsion-freeness of the
Levi-Civita connection reads as
$\nabla_{\xi_1}\xi_2-\nabla_{\xi_2}\xi_1=[\xi_1,\xi_2]$, which completes the proof of
the formula for $R^{\Cal A}$. In particular, $R^{\Cal A}$ vanishes if and only if $R$
vanishes which happens if and only if $g$ is flat.

Proceeding to $\tn^{\Cal A}$, differentiating the second formula in \eqref{conn-def},
we conclude that to compute $\tn^{\Cal A}_{\xi_1}\tn^{\Cal
  A}_{\xi_2}\binom{\eta}{\Phi}$, we have to add
\begin{equation}\label{tn2}
\begin{pmatrix} R(\xi_2,\eta)(\xi_1) \\ -\nabla_{\xi_1}R(\xi_2,\eta) -
  R(\xi_1,\nabla_{\xi_2}\eta)+R(\xi_1,\Phi(\xi_2)) 
\end{pmatrix}
\end{equation}
  to the right hand side of \eqref{n2}. Likewise, to compute $\tn^{\Cal
    A}_{[\xi_1,\xi_2]}\binom{\eta}{\Phi}$, we have to subtract
  $R([\xi_1,\xi_2],\eta)$ from the bottom component of the right hand side of
  \eqref{brack-term}.

  Together with the description of $R^{\Cal A}$, this readily implies that the top
  component of $\widetilde{R}^{\Cal A}(\xi_1,\xi_2)\binom{\eta}{\Phi}$ is given by
  $$
  R(\xi_1,\xi_2)(\eta)+R(\xi_2,\eta)(\xi_1)-R(\xi_1,\eta)(\xi_2)
  $$
  which vanishes by the first Bianchi identity \eqref{Bianchi1}. In the bottom
  component, we have terms containing $\eta$ and terms containing $\Ph$ and we treat
  these terms separately. For the terms containing $\eta$, we can first use
  torsion-freeness of $\nabla$ to write the contribution from the bracket term as
  $$
  R([\xi_1,\xi_2],\eta)=R(\nabla_{\xi_1}\xi_2,\eta)-R(\nabla_{\xi_2}\xi_1,\eta). 
  $$
  Adding this to the contributions from \eqref{tn2}, we can use the Leibniz rule for
  $\nabla$ to write the result as
  $$
  -(\nabla_{\xi_1}R)(\xi_2,\eta)+(\nabla_{\xi_2}R)(\xi_1,\eta)=(\nabla_\eta
  R)(\xi_1,\xi_2). 
  $$
  Here we used the second Bianchi identity \eqref{Bianchi2} in the last step.

  To understand the terms involving $\Phi$, we have to interpret the term
  $R(\xi_1,\xi_2)\bullet\Ph\in\Ga(\frak o(TM))$ taking into account that the
  infinitesimal action of $\frak o(n)$ on itself is the adjoint action. Hence for
  $\zeta\in\frak X(M)$, we obtain
  $$
(R(\xi_1,\xi_2)\bullet\Ph)(\zeta)=R(\xi_1,\xi_2)(\Phi(\zeta))-\Ph(R(\xi_1,\xi_2)(\zeta)). 
  $$
  On the other hand, the additional terms from \eqref{tn2} map $\zeta$ to
  $$
  R(\xi_1,\Ph(\xi_2))(\zeta)+R(\Phi(\xi_1),\xi_2)(\zeta), 
  $$ and adding this to the above, we exactly obtain $-(\Phi\bullet
  R)(\xi_1,\xi_2,\zeta)$ and we obtain the claimed formula for $\widetilde{R}^{\Cal
    A}$. Of course, this vanishes if and only if $\nabla R=0$ and $\Phi\bullet R=0$
  for any $\Phi\in\Ga(\frak (TM))$. The second equation means that in each point
  $x\in M$, the value $R(x)\in\La^2T^*_xM\otimes\frak o(T_xM)$ lies in the kernel of
  the action of $\frak o(T_xM)$, which has a direct interpretation in terms of
  representation theory. Indeed, decomposing $R$ into the Weyl curvature $W$ and the
  Ricci tensor $\Ric$, it follows that $W(x)=0$ and $\Ric(x)=\la(x)g_x$. Hence
  $(M,g)$ has to be conformally flat and Einstein, which is well known to imply that
  $\lambda$ is constant and to be equivalent to $(M,g)$ having constant sectional
  curvature. Since constant sectional curvature also implies that $\nabla R=0$, this
  completes the proof.
\end{proof}

\end{document}
